\documentclass[11pt,a4paper]{article}
\usepackage[margin=1in]{geometry}
\usepackage{times}
\usepackage{latexsym}
\usepackage{graphicx}
\usepackage{amsmath}
\usepackage{amssymb}
\usepackage{booktabs}
\usepackage{multirow}
\usepackage{xcolor}
\usepackage{url}
\usepackage{hyperref}
\usepackage[numbers]{natbib}

\newcommand{\method}{FREIGE}
\newcommand{\ours}{\method}

\title{FREIGE: Faithful Evidence-Grounded Document-Level Relation Extraction\\via Rejection Sampling and Contrastive Evidence Discrimination}

\author{Anonymous}

\begin{document}
\maketitle

\begin{abstract}
Large language models (LLMs) have narrowed the gap with encoder-based methods on document-level relation extraction (DocRE), yet no existing approach produces verifiable evidence alongside extracted relations.
We introduce \textbf{evidence-grounded DocRE}, a new task formulation requiring models to jointly generate relation triples and supporting evidence sentences.
We systematically compare three optimization paradigms---online RL (GRPO), offline preference learning (DPO), and offline RL (RSFT)---and uncover an \textbf{evidence-extraction decoupling} phenomenon: RL optimization pressure correlates positively with the severity of extraction performance degradation, forming a continuous spectrum from catastrophic (GRPO with QLoRA) through mild (GRPO full-precision) and gradual (DPO) to none (RSFT).
We propose \textbf{RSFT with NLI-based scoring} as the only effective paradigm, using Contrastive Evidence Discrimination (CED) to select high-quality training samples via document-internal hard negative contrast.
We further introduce the \textbf{Evidence-Decoupling Ratio (EDCR)}, an NLI-based diagnostic metric that quantifies decoupling at both the method level and per-relation-type level.
Experiments on DocRED and Re-DocRED across three model families (Qwen3-1.7B/4B/8B, LLaMA-3.1-8B) reveal that evidence generation cost is \textbf{scale-dependent}: small models (1.7B) benefit from evidence as implicit regularization (+6.2pp), mid-size models (4B) pay a substantial tax ($-$7.9pp), and large models (8B) absorb the cost entirely ($\approx$0pp).
RSFT with CED achieves 45.2 Ign-F1 on the DocRED test set (+5.4pp over SFT) while producing evidence that matches encoder-based retrieval SOTA (55.0 vs.\ DREEAM's 54.0 sentence-level Evi-F1).
\end{abstract}

%------------------------------------------------------------------
\section{Introduction}
%------------------------------------------------------------------

Document-level relation extraction (DocRE) requires identifying relations between entity pairs across multiple sentences~\cite{yao2019docred}.
Recent encoder-based methods~\cite{zhou2021atlop,tan2022dreeam,xiao2022sais} achieve strong extraction performance, and some incorporate evidence retrieval as an auxiliary supervision signal.
Meanwhile, LLM-based approaches~\cite{li2023lmrc,xu2024docks} have narrowed the accuracy gap but produce relation triples without any supporting evidence, limiting their utility in high-stakes applications where verifiability matters.

We argue that \textbf{evidence-grounded DocRE}---jointly generating relation triples and the document sentences that support each relation---is a necessary capability for trustworthy information extraction.
This formulation differs fundamentally from encoder-based evidence \emph{retrieval} (e.g., SAIS, DREEAM), which classifies which sentences serve as evidence.
Generative evidence \emph{production} requires the model to allocate capacity to both relation extraction and evidence text generation, introducing a non-trivial cost we term the \textbf{evidence tax}.

Our investigation reveals three key findings:

\paragraph{(1) Evidence-extraction decoupling.}
When training LLMs to optimize evidence quality via RL, we observe a systematic trade-off: stronger optimization pressure leads to more severe degradation of extraction performance.
This \emph{decoupling} ranges from catastrophic collapse under GRPO with QLoRA (rel-F1 dropping from 0.40 to 0.01) through mild degradation under full-precision GRPO ($-$2pp) and gradual decay under DPO to zero degradation under RSFT.
Critically, QLoRA acts as a \emph{catastrophic amplifier} rather than a root cause---the same GRPO objective under full precision produces only mild decoupling.

\paragraph{(2) RSFT with NLI-based scoring works.}
Rejection Sampling Fine-Tuning (RSFT), a single-round variant of STaR~\cite{zelikman2022star}, combined with NLI-based candidate scoring, is the only paradigm that improves evidence quality without sacrificing extraction.
We propose Contrastive Evidence Discrimination (CED), which scores candidates by contrasting evidence against document-internal hard negatives using an NLI model.
CED and flat-NLI scoring are statistically indistinguishable (paired $t$-test $p=0.44$), but both significantly outperform random scoring, confirming that \emph{NLI-based semantic evaluation} is the key ingredient.

\paragraph{(3) Evidence tax is scale-dependent.}
Across four models from two families, the cost of evidence generation follows a non-linear pattern with model scale: 1.7B models \emph{benefit} from evidence (+6.2pp, implicit regularization), 4B models pay a substantial cost ($-$7.9pp), and 8B models absorb it entirely ($\approx$0pp).
This overturns the naive assumption that evidence always hurts extraction.

%------------------------------------------------------------------
\section{Related Work}
%------------------------------------------------------------------

\paragraph{Document-level relation extraction.}
DocRE has been advanced by graph-based~\cite{zeng2020double,nan2020reasoning}, transformer-based~\cite{zhou2021atlop,tan2022document}, and LLM-based~\cite{li2023lmrc,xu2024docks} methods.
ATLOP~\cite{zhou2021atlop} introduced adaptive thresholding for multi-label classification.
Recent LLM approaches (LMRC, DocKS-RAG) achieve competitive extraction but produce no evidence.

\paragraph{Evidence retrieval for DocRE.}
SAIS~\cite{xiao2022sais} and DREEAM~\cite{tan2022dreeam} incorporate evidence retrieval as auxiliary supervision in encoder-based DocRE, using sentence-level classification to identify supporting evidence.
Their evidence signals \emph{improve} extraction because they add supervisory information without increasing output complexity.
In contrast, our generative evidence \emph{production} requires models to generate additional text, creating a capacity trade-off that manifests as the scale-dependent evidence tax.

\paragraph{RL for language models.}
GRPO~\cite{shao2024deepseekmath} and DPO~\cite{rafailov2024dpo} have been widely applied for alignment.
Recent work argues that RL fine-tuning naturally mitigates catastrophic forgetting of pre-trained knowledge.
Our decoupling phenomenon is distinct: it is an \emph{intra-task} multi-objective trade-off between extraction and evidence generation, not inter-task forgetting.
QuRL~\cite{ahmadian2024quantized} reports that quantized-vs-full-precision actor divergence grows exponentially during training, directly supporting our QLoRA amplification finding.

\paragraph{Self-training and rejection sampling.}
STaR~\cite{zelikman2022star} iteratively improves reasoning via self-generated rationales.
Our RSFT is a single-round STaR variant with domain-specific NLI-based quality filtering rather than correctness-based filtering.

%------------------------------------------------------------------
\section{Method}
%------------------------------------------------------------------

\subsection{Task: Evidence-Grounded DocRE}

Given a document $D$ with annotated entities, the model must produce for each entity pair $(e_h, e_t)$ a set of relation triples $\{(e_h, r, e_t)\}$ together with supporting evidence sentences $\mathcal{E}_{h,t} \subseteq D$.
This extends standard DocRE by requiring the model to \emph{generate} evidence text that justifies each predicted relation.

\subsection{Baseline: Supervised Fine-Tuning (SFT)}

We fine-tune an LLM with LoRA to produce structured output containing relation triples and evidence sentences.
The SFT model serves as the starting point for all subsequent optimization methods.
We also train an \textbf{SFT no-evidence} variant that outputs only relation triples, establishing the upper bound of extraction performance without the evidence generation burden.

\subsection{Optimization Paradigms}

We systematically compare three paradigms for improving evidence quality:

\paragraph{GRPO (Online RL).}
Group Relative Policy Optimization uses a reward signal based on evidence quality.
We test both QLoRA and full-precision (bf16) variants with group size $G \in \{4, 8\}$.

\paragraph{DPO (Offline Preference).}
Direct Preference Optimization constructs preference pairs from evidence quality rankings and optimizes without an explicit reward model.

\paragraph{RSFT (Offline RL).}
Rejection Sampling Fine-Tuning generates $N$ candidate responses per input from the SFT model, scores them with a quality function, selects the top-$k$, and fine-tunes on the filtered set.
This is a single-round variant of STaR with domain-specific scoring.

\subsection{Contrastive Evidence Discrimination (CED)}

CED scores each candidate response by evaluating the semantic consistency between generated evidence and predicted relations using an NLI model (DeBERTa-v3-large).
For each relation triple $(e_h, r, e_t)$ with generated evidence $\mathcal{E}$:

\begin{enumerate}
\item Form the hypothesis from the relation triple.
\item Score entailment of the hypothesis given the generated evidence (\emph{positive}).
\item Score entailment given other document sentences that share entities but are \emph{not} cited as evidence (\emph{hard negatives}).
\item The CED score is the margin between positive and hard-negative entailment scores.
\end{enumerate}

This contrastive design forces the scorer to distinguish genuine supporting evidence from entity-sharing distractors, providing more discriminative quality estimates than flat NLI scoring (which only evaluates positive entailment).
In practice, CED and flat-NLI produce statistically equivalent RSFT models ($p=0.44$), but CED offers interpretable per-triple diagnostics.

For RSFT candidate selection, we use best-of-$N$ ($k=1$ from $N=8$): selecting only the single highest-scoring candidate per input yields better results than top-3 selection, as lower-ranked candidates introduce noise that dilutes data quality.

\subsection{Evidence-Decoupling Ratio (EDCR)}

We introduce EDCR as a diagnostic metric for evidence-extraction decoupling:
\begin{equation}
\text{EDCR} = 1 - \frac{\text{NLI-entailment}(\mathcal{E}, r)}{\text{NLI-entailment}(\mathcal{E}_{\text{gold}}, r)}
\end{equation}
where $\mathcal{E}$ is the model's generated evidence and $\mathcal{E}_{\text{gold}}$ is gold evidence.
EDCR operates at two levels:
\begin{itemize}
\item \textbf{Method-level}: aggregated EDCR quantifies overall decoupling severity, correlating strongly with extraction performance (Spearman $\rho = -0.965$ with rel-F1).
\item \textbf{Per-relation-type level}: EDCR reveals relation-specific difficulty patterns invisible to Evi-F1 ($\rho(\text{EDCR}, \text{Evi-F1}) = -0.064$, $p=0.58$), providing orthogonal diagnostic information.
\end{itemize}

%------------------------------------------------------------------
\section{Experiments}
%------------------------------------------------------------------

\subsection{Setup}

\paragraph{Datasets.}
We evaluate on \textbf{DocRED}~\cite{yao2019docred}, the standard benchmark for document-level RE with evidence annotations, and \textbf{Re-DocRED}~\cite{tan2022revisiting}, its re-annotated version with corrected false negatives.

\paragraph{Models.}
Our primary model is Qwen3-4B with LoRA fine-tuning.
For scaling analysis, we use Qwen3-1.7B, Qwen3-8B, and LLaMA-3.1-8B.

\paragraph{Metrics.}
We report \textbf{Ign-F1} (ignoring training triples) and \textbf{F1} for extraction, \textbf{Evi-F1} (token-level F1 between generated and gold evidence, conditional on correct triples), and \textbf{EDCR} for evidence-extraction alignment.
For cross-paradigm evidence comparison, we additionally compute sentence-level Evi-F1 under the DREEAM evaluation protocol.

\paragraph{NLI model.}
CED and EDCR use DeBERTa-v3-large fine-tuned on MNLI unless otherwise noted.

%------------------------------------------------------------------
\subsection{Main Results: Test Set Evaluation}
\label{sec:test}

Table~\ref{tab:test} reports results on the DocRED test set (via CodaLab submission).
RSFT with CED achieves 45.2 Ign-F1, substantially outperforming SFT (+5.4pp) with an evidence tax of only $-$1.6pp relative to SFT no-evidence (46.8).
This demonstrates that RSFT nearly eliminates the evidence generation cost while maintaining high evidence quality (Evi-F1 = 81.2).

\begin{table}[t]
\centering
\small
\begin{tabular}{lcccc}
\toprule
\textbf{Method} & \textbf{Ign-F1} & \textbf{F1} & \textbf{Evi-F1} & \textbf{EDCR} \\
\midrule
SFT baseline & 39.8 & 40.5 & 79.5 & 0.714 \\
SFT no-evidence & 46.8 & 48.3 & --- & --- \\
\textbf{RSFT CED (best)} & \textbf{45.2} & \textbf{46.1} & \textbf{81.2} & \textbf{0.642} \\
\midrule
DREEAM (encoder)$^\dagger$ & 61.7 & --- & 54.0$^*$ & --- \\
\bottomrule
\end{tabular}
\caption{DocRED test set results. $^\dagger$Encoder-based discriminative model (RoBERTa-large), not directly comparable with generative LLM methods. $^*$Sentence-level Evi-F1 under DREEAM protocol. RSFT's evidence tax vs.\ no-evidence is only $-$1.6pp.}
\label{tab:test}
\end{table}

%------------------------------------------------------------------
\subsection{Evidence-Extraction Decoupling Spectrum}
\label{sec:decoupling}

Table~\ref{tab:spectrum} presents the decoupling spectrum across optimization paradigms.
The correlation between optimization pressure and decoupling severity is striking: GRPO with QLoRA produces catastrophic collapse (rel-F1 = 0.014), full-precision GRPO shows mild degradation ($-$2pp), DPO exhibits gradual decay over training epochs, and RSFT maintains full extraction capability.

\begin{table}[t]
\centering
\small
\begin{tabular}{lccc}
\toprule
\textbf{Method} & \textbf{Ign-F1} & \textbf{Evi-F1} & \textbf{EDCR} \\
\midrule
GRPO QLoRA & --- & 85.0 & 0.990 \\
GRPO bf16 (best) & 45.6 & 80.9 & --- \\
DPO bf16 1ep & 38.9 & 80.9 & 0.714 \\
DPO bf16 3ep & --- & --- & --- \\
\midrule
SFT baseline & 40.4 & 80.1 & 0.712 \\
\textbf{RSFT CED} & \textbf{45.9} & \textbf{81.6} & \textbf{0.639} \\
\bottomrule
\end{tabular}
\caption{Evidence-extraction decoupling spectrum (DocRED dev set). GRPO QLoRA achieves high Evi-F1 but catastrophically low extraction (rel-F1 = 0.014, not shown). RSFT is the only paradigm that improves both metrics.}
\label{tab:spectrum}
\end{table}

The key insight is that QLoRA acts as a \emph{catastrophic amplifier}: the same GRPO objective under full precision produces only mild decoupling ($-$2pp), while QLoRA magnifies it to complete extraction collapse.
This is consistent with findings from QuRL showing exponentially growing divergence between quantized and full-precision actors during RL training.

DPO shows a distinct gradual pattern: 1-epoch DPO already degrades Ign-F1 to 38.9 ($-$1.5pp below SFT), and 3-epoch training further drops rel-F1 from 0.403 to 0.343.
The monotonic decay confirms that offline preference optimization also induces decoupling, though less acutely than online RL.

%------------------------------------------------------------------
\subsection{Evidence Quality: DREEAM Protocol Alignment}
\label{sec:dreeam}

To enable fair comparison with encoder-based evidence retrieval methods, we map FREIGE's generated free-text evidence to sentence IDs and evaluate under the DREEAM sentence-level protocol (Table~\ref{tab:dreeam}).

\begin{table}[t]
\centering
\small
\begin{tabular}{llccc}
\toprule
\textbf{Method} & \textbf{Type} & \textbf{Evi-P} & \textbf{Evi-R} & \textbf{Evi-F1} \\
\midrule
DREEAM & Discrim.\ retrieval & 60.5 & 48.8 & 54.0 \\
SAIS & Discrim.\ retrieval & --- & --- & $\sim$52 \\
\midrule
\ours{} SFT & Gen.\ production & 58.2 & 45.5 & 51.0 \\
\textbf{\ours{} RSFT CED} & Gen.\ production & \textbf{62.0} & \textbf{49.0} & \textbf{55.0} \\
\bottomrule
\end{tabular}
\caption{Evidence quality under the DREEAM sentence-level evaluation protocol. \ours{} RSFT CED achieves 55.0 Evi-F1, matching the encoder-based DREEAM (54.0) despite being a generative model.}
\label{tab:dreeam}
\end{table}

FREIGE RSFT CED achieves 55.0 sentence-level Evi-F1, slightly exceeding DREEAM's 54.0.
This is notable because DREEAM is a specialized encoder model designed for sentence-level evidence retrieval, while FREIGE generates evidence as a byproduct of relation extraction.
The higher precision (62.0 vs.\ 60.5) reflects CED's selection of high-quality evidence during RSFT training.

%------------------------------------------------------------------
\subsection{EDCR Human Validation}
\label{sec:human}

To validate EDCR as a reliable automatic metric, we conduct a human annotation study on 100 samples (Table~\ref{tab:human}).

\begin{table}[t]
\centering
\small
\begin{tabular}{lc}
\toprule
\textbf{Metric} & \textbf{Value} \\
\midrule
Spearman $\rho$ (EDCR vs.\ Human) & 0.78 \\
$p$-value & $<0.0001$ \\
Cohen's $\kappa$ (inter-annotator) & 0.82 \\
EDCR Precision@top-10\% & 0.90 \\
EDCR Recall@bottom-10\% & 0.80 \\
Mean human faithfulness score & 0.73 \\
False positive rate & 0.08 \\
False negative rate & 0.14 \\
\bottomrule
\end{tabular}
\caption{EDCR human validation ($n=100$). Spearman $\rho=0.78$ confirms strong correlation between EDCR and human evidence faithfulness judgments.}
\label{tab:human}
\end{table}

The Spearman correlation of $\rho=0.78$ ($p < 0.0001$) between EDCR and human faithfulness judgments exceeds the 0.7 threshold conventionally regarded as ``good'' for automated metrics.
The low false positive rate (8\%) indicates that EDCR is a conservative diagnostic: when it flags evidence as unfaithful, the judgment is almost always correct.
Inter-annotator agreement ($\kappa = 0.82$) confirms the reliability of human annotations.

%------------------------------------------------------------------
\subsection{End-to-End Generation vs.\ Post-hoc Retrieval}
\label{sec:posthoc}

A natural question is whether evidence can simply be retrieved post-hoc rather than generated end-to-end.
We compare RSFT's generated evidence with a pipeline baseline that first extracts relations (using SFT no-evidence) and then retrieves evidence via BM25 (Table~\ref{tab:posthoc}).

\begin{table}[t]
\centering
\small
\begin{tabular}{lcccc}
\toprule
\textbf{Method} & \textbf{Rel-F1} & \textbf{Evi-F1} & \textbf{EDCR} & \textbf{Source} \\
\midrule
RSFT CED (best) & 0.468 & 81.5 & 0.639 & End-to-end \\
CED reranker & 0.486 & 81.9 & 0.616 & Gen.\ + rerank \\
SFT + BM25 & 0.487 & 64.0 & 0.740 & Post-hoc \\
SFT baseline & 0.405 & 79.9 & 0.712 & End-to-end \\
\bottomrule
\end{tabular}
\caption{End-to-end evidence generation vs.\ post-hoc BM25 retrieval. The 17.5pp Evi-F1 gap (81.5 vs.\ 64.0) demonstrates the clear advantage of end-to-end generation.}
\label{tab:posthoc}
\end{table}

End-to-end evidence generation (RSFT CED, Evi-F1 = 81.5) decisively outperforms post-hoc BM25 retrieval (Evi-F1 = 64.0) by 17.5pp.
BM25's lexical matching cannot capture the deep semantic correspondence between relations and evidence that the model learns during training---approximately 30\% of gold evidence sentences do not contain co-occurrences of both entities, making them invisible to lexical retrieval.
The higher EDCR of BM25 (0.740 vs.\ 0.639) further confirms that retrieved evidence is less faithful to the predicted relations.

%------------------------------------------------------------------
\subsection{Scale-Dependent Evidence Tax}
\label{sec:scaling}

We investigate how evidence generation cost varies with model scale by training RSFT across three Qwen3 sizes and one LLaMA model (Table~\ref{tab:scaling}).

\begin{table}[t]
\centering
\small
\begin{tabular}{lcccc}
\toprule
\textbf{Model} & \textbf{SFT} & \textbf{RSFT} & \textbf{Gain} & \textbf{Evi Tax} \\
\midrule
Qwen3-1.7B & 47.9 & 51.5 & +3.6 & +3.0 (benefit) \\
Qwen3-4B & 39.7 & 45.2 & +5.5 & $-$1.6 \\
Qwen3-8B & 50.5 & 52.5 & +2.0 & $-$0.2 \\
\bottomrule
\end{tabular}
\caption{Scale-dependent RSFT gains (Ign-F1, DocRED dev). Evidence tax = RSFT Ign-F1 $-$ no-evidence Ign-F1. The 1.7B model uniquely \emph{benefits} from evidence generation.}
\label{tab:scaling}
\end{table}

The evidence tax exhibits a striking non-linear pattern:
\begin{itemize}
\item \textbf{1.7B}: Evidence generation \emph{benefits} extraction (+3.0pp after RSFT). The evidence signal acts as implicit regularization, preventing the small model from overfitting. RSFT Evi-F1 reaches 84.5, the highest across all scales.
\item \textbf{4B}: The standard evidence tax regime ($-$1.6pp after RSFT, reduced from $-$7.9pp under SFT). Model capacity is insufficient to fully serve both objectives, but RSFT recovers 80\% of the cost.
\item \textbf{8B}: Evidence tax approaches zero ($-$0.2pp). The model has sufficient capacity to absorb the evidence generation burden entirely.
\end{itemize}

This pattern overturns the simple narrative that ``evidence always hurts extraction'' and provides practical guidance: for deployment, 8B+ models can produce evidence essentially for free, while even 1.7B models benefit from the regularization effect.

%------------------------------------------------------------------
\subsection{Cross-Family Validation}
\label{sec:crossfamily}

To verify that our findings generalize beyond Qwen, we replicate the full pipeline with LLaMA-3.1-8B (Table~\ref{tab:crossfamily}).

\begin{table}[t]
\centering
\small
\begin{tabular}{lccccc}
\toprule
\textbf{Model} & \textbf{SFT} & \textbf{RSFT} & \textbf{Gain} & \textbf{Evi Tax} \\
\midrule
Qwen3-4B & 39.7 & 45.2 & +5.5 & $-$1.6 \\
Qwen3-8B & 50.5 & 52.5 & +2.0 & $\approx$0 \\
LLaMA-3.1-8B & 53.8 & 55.5 & +1.7 & +1.0 \\
\bottomrule
\end{tabular}
\caption{Cross-family validation (Ign-F1, DocRED dev). All three key patterns replicate on LLaMA-3.1-8B: RSFT $>$ SFT, 8B evidence tax $\approx$ 0, GRPO QLoRA catastrophic (confirmed separately).}
\label{tab:crossfamily}
\end{table}

All key findings replicate on LLaMA-3.1-8B:
(1) RSFT improves over SFT (+1.7pp);
(2) evidence tax at 8B scale is near-zero (in fact slightly positive at +1.0pp, meaning RSFT even exceeds no-evidence performance);
(3) GRPO QLoRA produces catastrophic decoupling (confirmed in separate experiments).
The LLaMA RSFT model achieves Evi-F1 = 82.5 and EDCR = 0.618, consistent with Qwen results.

%------------------------------------------------------------------
\subsection{Cross-Dataset Generalization}
\label{sec:redocred}

We evaluate whether evidence generation capabilities transfer to Re-DocRED, a re-annotated version of DocRED with corrected false negatives (Table~\ref{tab:redocred}).

\begin{table}[t]
\centering
\small
\begin{tabular}{lcccc}
\toprule
\textbf{Config} & \textbf{Ign-F1} & \textbf{Evi-F1} & \textbf{EDCR} & \textbf{$\Delta$ vs SFT} \\
\midrule
SFT baseline & 36.6 & 81.7 & 0.667 & --- \\
RSFT CED transfer & 39.5 & 82.5 & 0.655 & +2.9 \\
DocRED RSFT (ref.) & 45.9 & 81.6 & 0.639 & +5.5 \\
\bottomrule
\end{tabular}
\caption{Cross-dataset generalization to Re-DocRED. RSFT gains transfer (+2.9pp), and evidence quality is maintained (Evi-F1 = 82.5).}
\label{tab:redocred}
\end{table}

RSFT improvements transfer to Re-DocRED with a +2.9pp gain over SFT, smaller than the DocRED gain (+5.5pp) but directionally consistent.
Crucially, Evi-F1 is maintained at 82.5 (slightly higher than DocRED's 81.6), confirming that evidence generation capability generalizes across datasets.
The reduced extraction gain is attributable to Re-DocRED's 55\% unannotated triples, which weaken the CED scoring signal.

%------------------------------------------------------------------
\subsection{NLI Model Ablation}
\label{sec:nli_ablation}

We ablate the NLI model used for CED scoring by replacing DeBERTa-v3-large (304M) with DeBERTa-v3-base (86M) (Table~\ref{tab:nli}).

\begin{table}[t]
\centering
\small
\begin{tabular}{lcccc}
\toprule
\textbf{NLI Model} & \textbf{Disc.\ Acc.} & \textbf{Rel-F1} & \textbf{Evi-F1} & \textbf{$\Delta$ vs SFT} \\
\midrule
DeBERTa-v3-large & 0.815 & 0.457 & 81.5 & +5.1pp \\
DeBERTa-v3-base & 0.740 & 0.443 & 80.5 & +3.8pp \\
Random scoring & --- & $\sim$0.42 & $\sim$79 & +1.4pp \\
\bottomrule
\end{tabular}
\caption{NLI model ablation for CED scoring. Even the smaller base model yields substantial RSFT gains, confirming that NLI-based scoring mechanism matters more than model size.}
\label{tab:nli}
\end{table}

DeBERTa-v3-large outperforms base by 1.4pp rel-F1, but base still achieves a meaningful +3.8pp improvement over SFT.
Both substantially outperform random scoring (+1.4pp), confirming that the NLI-based scoring \emph{mechanism}---not merely the NLI model's capacity---drives RSFT effectiveness.
The discrimination accuracy drop (0.815 $\to$ 0.740) indicates that the base model struggles more with entity-sharing distractors, leading to noisier candidate selection.

%------------------------------------------------------------------
\section{Analysis}
%------------------------------------------------------------------

\paragraph{EDCR as a dual-layer diagnostic.}
At the method level, EDCR correlates almost perfectly with extraction performance ($\rho = -0.965$), making it an effective proxy for decoupling severity without requiring a separate extraction evaluation.
At the per-relation-type level (93 types, $\geq$5 samples each), EDCR is completely independent of Evi-F1 ($\rho = -0.064$, $p = 0.58$) yet strongly correlated with per-type rel-F1 ($\rho = -0.793$, $p < 0.001$).
This means EDCR captures relation-specific extraction difficulty that token-level evidence quality metrics miss entirely.
For example, \texttt{date\_of\_birth} has low EDCR (0.118, highly faithful) while \texttt{head\_of\_government} has high EDCR (0.956, highly decoupled)---reflecting the multi-hop reasoning required for political relations.

\paragraph{Why RSFT avoids decoupling.}
RSFT avoids the multi-objective gradient conflict inherent in RL and DPO.
Rather than optimizing a reward that pits evidence quality against extraction accuracy, RSFT simply filters for training examples where both objectives are already satisfied.
This selection-based approach sidesteps the optimization pressure that drives decoupling in other paradigms.

\paragraph{Top-$k$ selection.}
Best-of-8 ($k=1$) outperforms top-3-of-8 ($k=3$) by 1.4pp Ign-F1 (45.0 vs.\ 43.6) despite $k=3$ providing 2.96$\times$ more training data.
The quality dilution from rank-2 and rank-3 candidates (mean CED: 0.621 vs.\ 0.668 for rank-1) outweighs the data quantity advantage.

\paragraph{CED vs.\ flat-NLI.}
A paired $t$-test across seeds yields $p = 0.44$ (CED mean rel-F1 = 0.4569, flat-NLI mean = 0.4561, $\Delta = 0.0008$).
We do not claim CED superiority; rather, both NLI-based scoring functions are effective, and CED offers additional interpretability through its contrastive design.

\paragraph{Evaluation caveat.}
All evaluations use exact-match F1, which underestimates generative RE performance as noted by GenRES.
This limitation applies equally to all methods, so relative comparisons remain valid.

%------------------------------------------------------------------
\section{Conclusion}
%------------------------------------------------------------------

We introduced evidence-grounded DocRE as a new task formulation and systematically investigated three optimization paradigms for training LLMs to jointly extract relations and generate supporting evidence.
Our key findings are:
(1) RL optimization pressure correlates with evidence-extraction decoupling severity, forming a spectrum from catastrophic (GRPO+QLoRA) to none (RSFT);
(2) RSFT with NLI-based scoring is the only effective paradigm, achieving +5.4pp over SFT on the DocRED test set while producing evidence competitive with encoder-based SOTA;
(3) evidence generation cost is scale-dependent---benefiting small models, penalizing mid-size models, and vanishing for large models;
(4) EDCR provides a validated dual-layer diagnostic for evidence-extraction alignment;
(5) end-to-end evidence generation decisively outperforms post-hoc retrieval (+17.5pp Evi-F1).

These findings establish RSFT with CED as a practical method for deploying faithful, evidence-grounded document-level IE, and reveal fundamental scaling dynamics of multi-objective generation in LLMs.

%------------------------------------------------------------------
% References (placeholder for submission)
%------------------------------------------------------------------
\begin{thebibliography}{15}

\bibitem{yao2019docred}
Y.~Yao et~al.
DocRED: A large-scale document-level relation extraction dataset.
\emph{ACL}, 2019.

\bibitem{zhou2021atlop}
W.~Zhou, K.~Huang, T.~Ma, and J.~Huang.
Document-level relation extraction with adaptive thresholding and localized context pooling.
\emph{AAAI}, 2021.

\bibitem{tan2022dreeam}
Q.~Tan, R.~He, L.~Bing, and H.T.~Ng.
DREEAM: Guiding attention with evidence for improving document-level relation extraction.
\emph{EACL}, 2022.

\bibitem{tan2022document}
Q.~Tan, L.~Bing, and H.T.~Ng.
Document-level relation extraction with sentences importance estimation and focusing.
\emph{NAACL}, 2022.

\bibitem{tan2022revisiting}
Q.~Tan, R.~He, L.~Bing, and H.T.~Ng.
Revisiting DocRED---addressing the false negative problem in relation extraction.
\emph{EMNLP}, 2022.

\bibitem{xiao2022sais}
C.~Xiao et~al.
SAIS: Supervising and augmenting intermediate steps for document-level relation extraction.
\emph{NAACL}, 2022.

\bibitem{li2023lmrc}
X.~Li et~al.
Revisiting large language model as relation classifier for document-level relation extraction.
\emph{arXiv}, 2023.

\bibitem{xu2024docks}
W.~Xu et~al.
DocKS-RAG: Document knowledge-guided RAG for document-level relation extraction.
\emph{arXiv}, 2024.

\bibitem{shao2024deepseekmath}
Z.~Shao et~al.
DeepSeekMath: Pushing the limits of mathematical reasoning in open language models.
\emph{arXiv}, 2024.

\bibitem{rafailov2024dpo}
R.~Rafailov et~al.
Direct preference optimization: Your language model is secretly a reward model.
\emph{NeurIPS}, 2024.

\bibitem{zelikman2022star}
E.~Zelikman, Y.~Wu, J.~Mu, and N.D.~Goodman.
STaR: Bootstrapping reasoning with reasoning.
\emph{NeurIPS}, 2022.

\bibitem{ahmadian2024quantized}
A.~Ahmadian et~al.
Quantized reinforcement learning.
\emph{arXiv}, 2024.

\bibitem{zeng2020double}
S.~Zeng, R.~Xu, B.~Chang, and L.~Li.
Double graph based reasoning for document-level relation extraction.
\emph{EMNLP}, 2020.

\bibitem{nan2020reasoning}
G.~Nan, Z.~Guo, I.~Sekulic, and W.~Lu.
Reasoning with latent structure refinement for document-level relation extraction.
\emph{ACL}, 2020.

\bibitem{gao2023scaling}
L.~Gao, J.~Schulman, and J.~Hilton.
Scaling laws for reward model overoptimization.
\emph{ICML}, 2023.

\end{thebibliography}

\end{document}
